\chapter{可用性与用户体验评价}\label{chap:evaluation}

设计层面的判断与算法层面的实现都需要回到「用户用得顺手吗」这一最朴素的检验。本章在小规模课堂环境下对 Mood Radio 进行了一轮形成性可用性评估——形成性意味着评估的目的不是给系统打一个最终分数，而是为下一轮迭代发现摩擦点。受课程时间所限，本评估的样本量很小（$N = 10$），所有结论应当被理解为「方向性」的；评估的真正价值在于揭示出哪些设计直觉得到了印证、哪些假设需要被重新审视。

\section{评估目标与方法}

围绕\cref{chap:requirement}~中提出的用户需求与设计原则，本次评估聚焦在三个相互关联的问题上：用户能否在不依赖任何引导的情况下完成「打开 $\to$ 触发推荐 $\to$ 播放」这条最关键的任务路径？用户能否解释系统是基于什么依据做出推荐？用户对系统在隐私、视觉反馈、可控性等关键属性上的表现是否满意？

为了让这三个问题各自得到回答，评估采用了三段式的混合方法。第一段是\textbf{有声思考任务}：受试者在不被引导的前提下依次完成 T1（摄像头识别 + 推荐 + 播放）、T3（在 \texttt{/chat} 用自然语言描述并播放）、T7（在 \texttt{/library} 中筛选 + 喜欢/不喜欢）三项任务，主持人记录完成情况、耗时与每一次明显的卡顿，并不在过程中给予任何提示。第二段是\textbf{5 级 Likert 满意度问卷}：任务结束后立即填写一份 7 题的量表问卷，覆盖任务难度、推荐理解度、隐私信任、视觉反馈、整体满意度等维度。最后第三段是\textbf{半结构化访谈}，主持人针对前两段中浮现出来的具体细节追问 5--10 分钟，让受试者补充未在问卷中体现的意见。

10 名同班同学接受了评估，男女比例约 6:4，年龄集中在 20--23 岁之间，其中 4 人首次使用本系统，6 人曾在原型阶段短暂接触过。每位受试者评估总时长约 25 分钟。

\section{任务完成情况}

三项任务的完成情况汇总在\cref{tab:task-completion}~中。T1 与 T7 的完成率均为 10/10，平均耗时也都在一分钟之内；T3 在网络层面遇到了一次失败——某位受试者评估时所连的 LLM 服务正巧出现短暂超时，导致这一次聊天没能拿到合法的 JSON 响应。剩下九位则都按预期完成了从「打字描述情绪」到「播放推荐歌单」的闭环。

\begin{table}[!htb]
  \centering
  \caption{有声思考任务的完成情况（$N=10$）}\label{tab:task-completion}
  \begin{tabular*}{\linewidth}{@{\extracolsep{\fill}}llccp{5cm}}
    \toprule
    \textbf{编号} & \textbf{任务} & \textbf{完成率} & \textbf{平均耗时} & \textbf{主要卡点} \\
    \midrule
    T1 & 摄像头识别 + 一键推荐 + 播放 & 10/10 & 38\,s & 摄像头授权弹窗位置容易被忽略 \\
    T3 & 在 /chat 用自然语言描述并播放 & 9/10 & 64\,s & 一例因 LLM 配置网络问题失败 \\
    T7 & 在 /library 完成「筛选 + 喜欢/不喜欢」 & 10/10 & 51\,s & 部分受试者一开始未意识到卡片可点开详情 \\
    \bottomrule
  \end{tabular*}
\end{table}

「主要卡点」一栏对设计调整最有指导意义。在 T1 中，多名受试者反映「不知道授权弹窗会出现在哪个位置」——这本质上是浏览器的平台问题（Chrome 的弹窗在地址栏附近，Edge 在更靠右的位置，火狐则会暗弹一个小气泡），并不能由 Mood Radio 直接控制；但课题组可以通过在按下「开启摄像头」之后先弹出一张「请允许浏览器访问摄像头」的引导提示来缓解这一困惑。在 T7 中，「卡片可以点开详情」是一个隐式的交互约定，初次使用的受试者并不知道这一点——这提示课题组在卡片右下角加一个小尺寸的「更多」按钮或者轻微的悬停动效，让可点击性更显眼。

\section{用户满意度问卷}

\cref{tab:likert}~给出 7 题 5 级 Likert 量表的统计结果。「均值」越接近 5 越正面，「肯定比例」表示给出 4 分或以上的受试者占比。问卷的总体结论可以这样概括：受试者最满意的是「本地优先」的隐私设计（Q4，均值 4.6）与「任务可独立完成」的可达性（Q1，均值 4.4）；相对来说，他们对「推荐理由的可解释性」（Q3，均值 4.0）与「向他人推荐意愿」（Q7，均值 4.0）的评分要克制一些。前者印证了课题组在「本地优先」与「低负担入口」这两条原则上的投入是值得的；后者则提示「推荐徽标」这一可解释手段虽然让大部分人理解了打分逻辑，但还未达到「我会主动安利给朋友」的程度——这与系统在曲库规模、推荐多样性方面仍有改善空间是一致的。

\begin{table}[!htb]
  \centering
  \caption{用户满意度问卷（5 级 Likert，$N=10$）}\label{tab:likert}
  \begin{tabularx}{\linewidth}{>{\ttfamily}lX>{\centering\arraybackslash}p{1.6cm}>{\centering\arraybackslash}p{1.6cm}}
    \toprule
    \textbf{编号} & \textbf{题项} & \textbf{均值} & \textbf{肯定比例} \\
    \midrule
    Q1 & 我可以在不看任何说明的情况下完成「触发推荐」 & 4.4 & 9/10 \\
    Q2 & 系统给出的推荐让我觉得「它理解我此刻的情绪」 & 4.1 & 8/10 \\
    Q3 & 我能理解为什么系统会推荐这些歌曲（可解释性） & 4.0 & 8/10 \\
    Q4 & 我相信系统的「本地优先」隐私设计 & 4.6 & 9/10 \\
    Q5 & 界面视觉变化（背景、主题、黑胶）能体现系统对情绪的回应 & 4.3 & 8/10 \\
    Q6 & 「匹配 / 安抚」双策略对我来说是有必要的 & 4.2 & 8/10 \\
    Q7 & 我愿意把这个系统介绍给身边的同学 & 4.0 & 7/10 \\
    \bottomrule
  \end{tabularx}
\end{table}

\section{开放反馈与问题分析}

定量结果之外，访谈环节贡献了若干更具体的意见，可以归结为五个具有共性的方向。\emph{首先}是\textbf{首次扫描音乐库较慢}的问题：当个人音乐库包含数百首歌时，初次扫描需要完成解码、Mel 频谱与情绪预测，过程耗时几十秒到几分钟不等，多位受试者表示「我以为系统卡死了」。这并非一个算法问题，更多是一个「状态告知不足」的问题——可以在扫描期间显式给出已完成数量、总数量以及预估剩余时长，让等待变得可被感知。\emph{其次}是\textbf{弱光下表情识别不稳定}：face-api.js 在背光、侧脸或遮挡场景中置信度会显著下降，「系统老识别成 neutral」是访谈中的高频说法；从产品角度看，与其试图改进识别本身，不如在置信度低于某个阈值时主动建议用户「不如手动选一下情绪」，这与「多通道冗余」的设计哲学也是一致的。\emph{第三}是\textbf{小型曲库下的推荐多样性不足}：当个人本地音乐库小于 50 首时，复合打分会反复推送几首高分曲目；改进方向是在打分中加入「同情绪曲目轮转」的项，避免高分曲目长期占据榜首。\emph{第四}是\textbf{聊天的偶发网络失败}：T3 中那次失败暴露了系统在错误状态下的处理缺失——当前页面没有清晰地提示「LLM 配置可能有问题，请检查 base\_url」，而只是停在了 loading 状态。这一问题虽小，却恰恰击中了 Nielsen 启发式中的「适当的错误恢复」。\emph{最后}是\textbf{授权弹窗位置}的混淆：这一问题严格来说是浏览器平台层的责任，但作为应用层我们仍可以在点击开启摄像头按钮之后先弹出一张引导图，告诉用户「请在地址栏左侧的弹窗中点 \emph{允许}」。

这些反馈大多并不要求大动干戈的重构，反而恰恰是「细节决定体验」的常见案例。课题组把它们列入了\cref{chap:conclusion}~讨论的后续工作。

\section{HCI 原则映射}

把 Mood Radio 的功能特性与人机交互课程中讨论的核心概念做一次对照，可以帮助我们检视系统的功能集合是否真的围绕一组明确的设计理念展开，而非散点功能的堆砌。\cref{tab:hci-mapping}~给出这一对照。

\begin{table}[!htb]
  \centering
  \caption{Mood Radio 与 HCI 课程概念的映射}\label{tab:hci-mapping}
  \begin{tabularx}{\linewidth}{lX}
    \toprule
    \textbf{HCI 概念} & \textbf{项目体现} \\
    \midrule
    多模态输入       & 摄像头表情 + 鼠标点击 + 自然语言 + 头部手势 + 时段 \\
    隐式 vs 显式输入 & 表情/时段为隐式，手动/聊天为显式 \\
    情境感知         & 时段电台、表情识别、心情日志均以「上下文」为信号 \\
    沉浸感设计       & Aurora 背景 + 黑胶旋转 + 字体配对 + 情绪驱动主题 \\
    反馈环路         & 推荐 $\to$ 用户喜欢/不喜欢 $\to$ 复合打分调整 $\to$ 下一次推荐 \\
    群体在场（presence） & 匿名 UUID + 心跳聚合 + 隐私门槛 \\
    可解释性         & \texttt{debug\_score} 徽标 + 聊天解释 \\
    无障碍           & Skip link + aria-pressed + 语义化 heading + 键盘导航 \\
    隐私优先         & 摄像头本地处理 + LLM key 不落盘 + 在场聚合不返回个人 \\
    \bottomrule
  \end{tabularx}
\end{table}

\section{在线 / 离线运行方式}

为方便课程评审和后续使用者复用，Mood Radio 提供「离线本地运行」与「仅前端调试运行」两种入口，两者都以 \texttt{start.ps1} 一键脚本作为统一入口，差别仅在传入参数。\cref{tab:run-mode}~给出三种参数组合的含义。除聊天功能依赖外部 LLM 服务之外，其余所有功能都完全在本机运行；后端默认监听 \texttt{127.0.0.1:8080}，前端 dev server 默认监听 \texttt{5173}，并通过 Vite 的 \texttt{/api $\to$ 8080} 反向代理对接。

\begin{table}[!htb]
  \centering
  \caption{Mood Radio 的运行方式}\label{tab:run-mode}
  \begin{tabularx}{\linewidth}{lXX}
    \toprule
    \textbf{模式} & \textbf{命令} & \textbf{说明} \\
    \midrule
    完整离线运行 & \texttt{.{\textbackslash}start.ps1} & 后端新窗口 + 前端 dev；除聊天功能外完全离线 \\
    仅后端       & \texttt{.{\textbackslash}start.ps1 -BackendOnly}  & 只跑 C++ 后端 + SQLite，便于联调推荐接口 \\
    仅前端       & \texttt{.{\textbackslash}start.ps1 -FrontendOnly} & 只跑 Vite dev server，便于 UI 改样式时热刷 \\
    \bottomrule
  \end{tabularx}
\end{table}

\cref{lst:run}~给出一次性的安装命令与启动命令。聊天功能需要在 \texttt{/settings} 页面填入 \texttt{base\_url}、\texttt{api\_key} 与 \texttt{model} 这三项 LLM 配置，这些参数仅存于浏览器 \texttt{localStorage}，不进入后端数据库，符合\cref{tab:nfr}~中的 NFR-02。

\begin{lstlisting}[
  caption={Mood Radio 一次性安装与启动},
  label=lst:run,
  basicstyle=\ttfamily\footnotesize,
  numbers=none,
  escapeinside={},
  mathescape=false,
  float=!htb,
]
# 一次性安装（仅首次需要）
cd frontend\web && npm install  ; cd ..\..
cd scripts && uv venv --python 3.12 .venv ; ^
  uv pip install -p .venv/Scripts/python.exe ^
  onnxruntime librosa soundfile numpy ; cd ..

# 编译后端（MSVC 工具链 + vcpkg）
cmake -B build -S . -DCMAKE_TOOLCHAIN_FILE=...
cmake --build build

# 启动（从仓库根）
.\start.ps1                # 后端新窗口 + 前端 dev
.\start.ps1 -BackendOnly   # 只跑后端
.\start.ps1 -FrontendOnly  # 只跑前端
\end{lstlisting}
